\documentclass[11pt,a4paper]{article}
\usepackage[margin=26mm]{geometry}
\usepackage{amsmath,amssymb,amsthm,booktabs,array}
\usepackage[T1]{fontenc}
\usepackage[colorlinks=true,allcolors=blue]{hyperref}
\usepackage{microtype}
\hypersetup{pdftitle={The section 2.9 conjecture of Lester, Haddadin and Gripaios (arXiv:2008.05206v6): three parity-odd variables are necessary},
pdfauthor={Haobo Ma; Wenlin Zhang},pdfsubject={Minimal continuous invariant parity-odd event variables}}
\newtheorem{theorem}{Theorem}
\newtheorem{lemma}[theorem]{Lemma}
\newtheorem{corollary}[theorem]{Corollary}
\newcommand{\RR}{\mathbb R}
\newcommand{\SO}{\mathrm{SO}^{+}(1,3)}
\newcommand{\PP}{\mathcal P}
\newcommand{\all}{\Omega_{\mathrm{all}}}
\newcommand{\coll}{\Omega_{\mathrm{coll}}}
\title{The \S2.9 conjecture of Lester, Haddadin and Gripaios\break
(arXiv:2008.05206v6):\break
three parity-odd variables are necessary}
\author{Haobo Ma\\\small ChronoAI Pte Ltd, Singapore
\and Wenlin Zhang\\\small National University of Singapore\\
\small ChronoAI Pte Ltd, Singapore}
\date{3 November 2026}
\begin{document}
\maketitle
\begin{abstract}
We prove the conjecture in \S2.9 of Lester, Haddadin and Gripaios,
\emph{Lorentz and permutation invariants of particles III}
(arXiv:2008.05206v6), that their three-variable sets $S$ and $\hat S$
are minimal for their respective four-particle event classes.
Any two continuous, Lorentz- and pair-permutation-invariant, parity-odd
real variables vanish together on a chiral collision event.
The proof constructs a sphere consisting entirely of chiral events,
on which parity acts antipodally up to the allowed symmetries, and
applies the two-dimensional Borsuk--Ulam theorem.
The same obstruction holds for strictly timelike, momentum-completable
collisions with any number of additional interchangeable momenta,
and for fully interchangeable collections of at least four momenta.
\end{abstract}

\section{The conjecture and event domains}

Lester, Haddadin and Gripaios construct a sufficient triple $S$ for
four-particle events with two interchangeable pairs, and a simpler
triple $\hat S$ for the collision subclass \cite[\S5.1]{LHG}.
Their conjecture in \S2.9 is that neither triple can be replaced by a
one- or two-variable set retaining continuity, invariance, parity oddness
and sufficiency. Here sufficiency means that the variables have no
common zero on a chiral event. The missing ingredient is a lower bound
that applies to \emph{every} admissible choice of variables, rather
than only to subsets of the displayed triples.

Write $x\cdot y=x^0y^0-x^1y^1-x^2y^2-x^3y^3$ and $x^2=x\cdot x$.
Let $V_+=\{x\in\RR^4:x^0\geq0,\ x^2\geq0\}$, including the zero
vector. An event is an ordered quadruple $e=(p,q,a,b)$, with domains
\[
 \all=V_+^4,\qquad
 \coll=\{(p,q,a,b)\in\all:p^2q^2<(p\cdot q)^2\}.
\]
These are the source's $\Omega_{(pq)(ab)X}$ and
$\Omega_{(pq)\rightarrow(ab)X}$, respectively
\cite[\S3, definition of collision events and final two domain definitions]{LHG}.
The domains allow variable masses and unobserved additional momenta;
they do not impose $p+q=a+b$.

The group $G=\SO\times S_2(pq)\times S_2(ab)$ acts by a common
proper orthochronous Lorentz transformation and independent exchanges
within the two pairs. Parity is $\PP x=(x^0,-x^1,-x^2,-x^3)$,
applied to each momentum. An event is \emph{chiral} if $g e\ne\PP e$
for every $g\in G$.
For either domain $\Omega$, call $f:\Omega\to\RR$ admissible if it
is continuous in the subspace topology, $G$-invariant, and satisfies
$f(\PP e)=-f(e)$. Both domains are preserved by $G$ and parity.

\begin{theorem}\label{thm:main}
For $\Omega=\all$ or $\Omega=\coll$, and any two admissible variables
$f_1,f_2$, there is a chiral event $e\in\Omega$ with
$f_1(e)=f_2(e)=0$. Consequently a sufficient admissible set on either
domain has at least three variables. The sets $S$ and $\hat S$ of
\cite[\S5.1]{LHG} are therefore minimal on their respective domains.
\end{theorem}

\section{A sphere of chiral collision events}

For $v=(s,t,d)\in S^2$, so $s^2+t^2+d^2=1$, put
$A=\sqrt{4+t}$ and $B=\sqrt{4-t}$ and define
\begin{equation}\label{eq:family}
 \begin{aligned}
 p&=(30,0,0,1),& q&=(20,0,0,-1),\\
 a_v&=(10+s,A,0,d),& b_v&=(10-s,0,B,-d),\\
 e_v&=(p,q,a_v,b_v).
 \end{aligned}
\end{equation}
The map $v\mapsto e_v$ is continuous. Since $|s|,|t|,|d|\leq1$,
we have $A,B\geq\sqrt3$, $a_v^0,b_v^0\geq9$, and each of their
spatial squared norms is at most $6$. Thus all four momenta are
strictly future timelike. Moreover
\begin{equation}\label{eq:collision}
 p^2=899,\qquad q^2=399,\qquad p\cdot q=601,\qquad
 (p\cdot q)^2-p^2q^2=2500>0.
\end{equation}
In particular $e_v\in\coll$ for every $v$.

\begin{lemma}\label{lem:chiral}
Every event $e_v$ in \eqref{eq:family} is chiral.
\end{lemma}
\begin{proof}
Suppose a Lorentz transformation $\Lambda\in\SO$ and the allowed
exchanges carry $e_v$ to $\PP e_v$. An exchange of $p$ and $q$ is
impossible: Lorentz transformations and parity preserve squared norms,
but $p^2\ne q^2$. Hence the first pair is fixed and
$\Lambda p=\PP p$, $\Lambda q=\PP q$.

Set $H=\PP\Lambda$, viewing $\PP=\operatorname{diag}(1,-1,-1,-1)$
as a matrix. Then $H$ preserves the Minkowski form, has determinant
$-1$, and fixes $p$ and $q$. The two vectors span the time--longitudinal
plane: $p+q=50e_0$ and $p-30e_0=e_3$. Therefore $H$ fixes $e_0,e_3$
and preserves their Minkowski-orthogonal complement, the transverse
Euclidean plane. Its restriction $U$ to that plane belongs to $O(2)$
and has determinant $-1$.

If the second pair is also fixed, $H a_v=a_v$ and $H b_v=b_v$.
Their transverse components are $A e_1$ and $B e_2$, with $A,B>0$.
Thus $U e_1=e_1$ and $U e_2=e_2$, contradicting $\det U=-1$.

If the second pair is exchanged, $H a_v=b_v$. The fixed time
coordinate gives $10+s=10-s$, hence $s=0$. The fixed longitudinal
coordinate gives $d=-d$, hence $d=0$. Orthogonality of $U$ gives
$A^2=B^2$, so $4+t=4-t$ and $t=0$. This contradicts
$s^2+t^2+d^2=1$. All allowed exchanges have been exhausted.
\end{proof}

Let $R$ fix time and act on space by
\[
 (x,y,z)\longmapsto(-y,-x,-z).
\]
This spatial rotation has determinant $1$. Its time-fixing extension
belongs to $\SO$. Direct substitution gives
\begin{equation}\label{eq:antipodal}
 e_{-v}= R\,(ab)\,\PP e_v.
\end{equation}
For example, applying $R$ to the parity image of $b_v$ gives
$(10-s,B,0,-d)=a_{-v}$. The transformations of $p,q,a_v$ follow by
direct substitution. Consequently every admissible $f$ satisfies
\begin{equation}\label{eq:odd}
 f(e_{-v})=f(R(ab)\PP e_v)=f(\PP e_v)=-f(e_v).
\end{equation}

\section{The topological obstruction and minimality}

The Borsuk--Ulam theorem states that every continuous map
$S^2\to\RR^2$ takes equal values at some antipodal pair
\cite[Satz II, p.~178]{Borsuk}. For an odd map that value is zero.
For completeness, the required odd-map form also has the following
short proof. If an odd continuous map $F:S^2\to\RR^2$ had no zero,
normalization would give an odd map $g:S^2\to S^1$. On the closed
upper hemisphere, a disk, $g$ has a continuous real lift $\theta$ under
$u\mapsto e^{iu}$. On the equator,
$\theta(-z)-\theta(z)\in\pi+2\pi\mathbb Z$ by oddness.
This difference is continuous on the connected equator, hence constant.
Replacing $z$ by $-z$ negates it; the constant would be zero, impossible.

Apply this fact to
\[
 F(v)=\bigl(f_1(e_v),f_2(e_v)\bigr).
\]
Continuity follows from continuity of the two variables and the family,
and \eqref{eq:odd} makes $F$ odd. A zero supplies the desired event,
which is chiral by Lemma~\ref{lem:chiral} and belongs to the collision
domain by \eqref{eq:collision}. This proves the assertion for
$\coll$, and also for $\all$ by restricting the two variables to
$\coll$. A single variable is ruled out by taking $f_2=0$.
The source's sufficient triples then give the matching upper bounds
and prove Theorem~\ref{thm:main}.

Scalar continuity obstructions have an earlier counterpart in chiral
connectedness of Euclidean shapes, discussed by Fowler in the volume
dated 2005 \cite{Fowler}. Wang's March 2021 two-component tetrahedral
descriptor \cite{Wang} concerns Euclidean point configurations modulo
translations and rotations, rather than Lorentzian four-momentum events.
Lester's November 2021 \emph{Chiral Measurements} uses paths of chiral
objects to obstruct a single continuous ideal scalar parity
\cite[\S3.3, Lemma 3.8]{Lester}. The sphere above gives the stronger
two-component obstruction for these event classes. It is the antipodal
two-sphere, rather than a path alone, that forces the simultaneous zero
of two arbitrary variables.

\section{Consequences and coverage}

\begin{corollary}\label{cor:strict}
The two-variable obstruction remains valid on the subdomain of
strictly future-timelike collision events. It also remains valid when
an additional future-timelike unobserved momentum is required to
complete momentum conservation.
\end{corollary}
\begin{proof}
The family already lies in the strict timelike subdomain. In addition,
\[
 k_v=p+q-a_v-b_v=(30,-A,-B,0),\qquad k_v^2=900-(4+t)-(4-t)=892.
\]
Thus every event of the family has the requisite future-timelike
completion. These subdomains are parity- and $G$-invariant, so the
same continuous odd-map argument applies to variables on them.
\end{proof}

\begin{corollary}\label{cor:many}
For every integer $n\geq2$, on the full forward-cone domain of an
interchangeable pair and $n$ further interchangeable momenta, any
two continuous $\SO\times S_2\times S_n$-invariant parity-odd real
variables vanish together on a chiral event. The same holds on the
collision subclass defined by the first pair.
The conclusion also holds for strictly timelike collisions admitting
a future-timelike unobserved momentum that completes conservation.
\end{corollary}
\begin{proof}
Put $r=n-2$ and append $r$ identical momenta
$c_j=((r+1)^{-1},0,0,0)$ to $e_v$. Their squared masses are at most
$1$, whereas $75\leq a_v^2,b_v^2\leq121$. Thus any allowed
Lorentz/parity mapping preserves $\{a_v,b_v\}$ and the collection of
markers separately. Lemma~\ref{lem:chiral} rules out such a mapping.
Parity and $R$ fix the markers, so \eqref{eq:antipodal} still holds.
All momenta are strictly future timelike, and the unobserved completion is
\[
 k_{v,r}=\left(30-\frac{r}{r+1},-A,-B,0\right),\qquad
 k_{v,r}^2=833+\frac{58r+59}{(r+1)^2}>833.
\]
Consequently the assertion also holds on the subdomain of strictly
timelike collisions admitting a future-timelike unobserved completion.
These domains are symmetry invariant, so Borsuk--Ulam applies as before.
\end{proof}

\newpage
\begin{corollary}\label{cor:permutations}
Let $N\geq4$ and let $H\leq S_N$ contain the transposition of the
third and fourth slots. On $V_+^N$, any two continuous
$\SO\times H$-invariant parity-odd real variables have a common zero
at an $H$-chiral event. This holds even when all momenta are required
to be strictly future timelike. In particular it holds for $H=S_N$,
and for $S_m\times S_n$ with $m,n\geq2$ after relabeling the slots.
\end{corollary}
\begin{proof}
Use $e_v$ with $r=N-4$ markers from Corollary~\ref{cor:many}.
Under any Lorentz/permutation mapping, $p$ and $q$ are individually
fixed by their unique squared masses $899$ and $399$; the pair
$\{a_v,b_v\}$ is separated from the markers by squared mass.
The chirality argument therefore works even for the full group $S_N$.
The permitted transposition supplies \eqref{eq:antipodal}, giving the
same odd two-sphere. For two blocks, place $p,q$ in the first block,
$a,b$ in the second, and the remaining markers in either block.
\end{proof}

\begin{corollary}\label{cor:scalar}
On either four-momentum domain, a continuous invariant parity-odd
complex-valued scalar also vanishes on some chiral event. In particular,
neither a real nor a complex continuous scalar can be sufficient.
If continuity is dropped, one real variable is sufficient.
\end{corollary}
\begin{proof}
Apply Theorem~\ref{thm:main} to the real and imaginary parts; a real
scalar is the special case with zero imaginary part. Conversely, for
either sufficient triple $(X_1,X_2,X_3)$, return its first nonzero
component, or zero if all vanish. This is invariant and parity odd,
since parity preserves the index of the first nonzero component.
Its zero set is exactly the achiral set. This is the construction in
\cite[\S6.1.1, (226)]{LHG}; Theorem~\ref{thm:main} proves that its
restriction to the chiral sphere must be discontinuous.
\end{proof}

\begin{corollary}\label{cor:labelled}
For four labeled forward-causal momenta, with no permutation symmetry,
the single continuous variable $D(e)=\det[p\ q\ a\ b]$ is sufficient.
For at most three forward-causal momenta every event is achiral.
\end{corollary}
\begin{proof}
The determinant is Lorentz invariant and parity odd. If it is nonzero,
no proper Lorentz transformation maps the labeled event to its parity
image. If the span $W$ of the momenta has dimension at most three and
contains a timelike vector, $W^\perp$ contains a spacelike vector $u$.
Reflection in $u^\perp$ fixes $W$, has determinant $-1$, and is
orthochronous because it fixes a timelike vector. Composing it with
parity gives a proper orthochronous transformation taking the event to
its parity image. If $W$ contains no timelike vector, all nonzero
forward-causal momenta lie on a single null ray: the sum of two
nonparallel future null vectors is timelike. A purely spatial vector
orthogonal to that ray supplies the same reflection; if all momenta
vanish choose any spatial axis. Thus $D=0$ implies achirality for four
labeled momenta, and every event with at most three momenta is achiral.
\end{proof}

\clearpage
\paragraph{Coverage of source questions.}
The classifications below distinguish a proved lower bound from a
determination of the minimum, and from independence of a displayed set.
\smallskip

\noindent\begin{tabular}{@{}>{\raggedright\arraybackslash}p{.40\linewidth}
>{\raggedright\arraybackslash}p{.55\linewidth}@{}}
\toprule
Source item in arXiv:2008.05206v6 & Coverage\\
\midrule
\S2.9: minimality of $S$ and $\hat S$ & Proved: minimum three,
Theorem~\ref{thm:main} and cited \S5.1 upper bounds.\\[3pt]
\S6.1, first skeptic's scalar alternative & Proved impossible on these
domains, even for a complex scalar (Corollary~\ref{cor:scalar}).\\[3pt]
\S6.1.1, discontinuous compression & Prior proof cited: (226),
reproduced in Corollary~\ref{cor:scalar}.\\[3pt]
\S2.5, $S_2\times S_n$ extension & Proved: lower bound three for all
$n\geq2$, also with timelike completion (Corollary~\ref{cor:many}).
Construction of efficient variables and empirical performance: not covered.\\[3pt]
\S5.2.8, $S_{19}$ minimality & Proved: lower bound three. Exact minimum,
lower-degree alternatives and exploiting $C_{4B}$: not covered.\\[3pt]
\S5.2.6, parameter-dependent irreducibility & Not covered: necessity
of all nineteen variables for arbitrary nonzero $u_i$.\\[3pt]
\S5.2.7, fixed-mass $S_9$ & Not covered: equal-mass incoming beams
and three massless outgoing momenta.\\[3pt]
\S5.4, $S_{28}$ & Proved: lower bound three. Exact minimum and
independence of the twenty-eight displayed variables: not covered.\\[3pt]
\S6.2, restriction $c\to0$ & Prior proof cited: restriction of the
five-momentum formulas. No new claim about their independence.\\
\bottomrule
\end{tabular}

\smallskip
For the cross-referenced \emph{Chiral Measurements}, \S3.3,
Lemma~3.8 gives a prior scalar path obstruction; the present two-variable
obstruction strengthens it on the specified event domains.
The existence speculation in its \S3.4 is realized here by the cited
triples, but termination of the proposed general algorithm is not covered.
Its \S2.8 continuity conjecture for general chiral measures is also
not covered. Exact minima for the five-momentum classes and fixed-mass
variants remain separate questions.

\section*{Mathematical verification and reproducibility}

The Lean declaration
\nolinkurl{D5.S3.Geometry.Lorentz.ChiralEventParityMinimality.result}
in the pinned
\href{https://github.com/the-omega-institute/trureturing/blob/50f2a97bb5e704842f2c32c6095f68a0dc62c779/D5/S3/Geometry/Lorentz/ChiralEventParityMinimality.lean}{\texttt{ChiralEventParityMinimality.lean}}
has the lower-bound statement of Theorem~\ref{thm:main} for
\texttt{allEvents} and \texttt{collisionEvents}.
Its \texttt{admissible}, \texttt{chiral} and \texttt{family}
definitions are the definitions and family used here.
The conclusion of exact minimality uses the cited sufficient triples.
Corollaries~\ref{cor:strict}--\ref{cor:labelled} have the written proofs
given above; their additional domains and permutation groups are not
conclusions of that Lean declaration. The script \texttt{verify.py}
checks the exact rotation, antipodal and collision identities, the
algebraic contradiction in the exchange case, the marker/completion
formulas and determinant identities. Its finite checks cover 31 rational
sphere inputs and the marker families with $4\leq N\leq12$.
The topological and reflection arguments establish the all-input results.

\clearpage
\section*{AI tools}
In preparing this manuscript, the authors used AI tools (ChatGPT Pro, Codex
and Claude) for literature search, computation, checking and formal
verification in Lean. All proofs can be checked by hand, and the
computations can be reproduced from the accompanying verification package.
The authors take full responsibility for the content.

\section*{License}
This note is licensed under
\href{https://creativecommons.org/licenses/by/4.0/}{Creative Commons
Attribution 4.0 International (CC BY 4.0)}.

\begin{thebibliography}{9}
\raggedright
\bibitem{LHG}
C.~G.~Lester, W.~Haddadin and B.~Gripaios,
\emph{Lorentz and permutation invariants of particles III:
constraining nonstandard sources of parity violation},
International Journal of Modern Physics A \textbf{37}(16) (2022), 2250093.
\href{https://doi.org/10.1142/S0217751X22500932}{doi:10.1142/S0217751X22500932}.
Posing text: \href{https://arxiv.org/abs/2008.05206v6}{arXiv:2008.05206v6},
\S2.9; sufficient triples: \S5.1.

\bibitem{Borsuk}
K.~Borsuk, \emph{Drei S\"atze \"uber die $n$-dimensionale euklidische
Sph\"are}, Fundamenta Mathematicae \textbf{20} (1933), 177--190.
\href{https://doi.org/10.4064/fm-20-1-177-190}{doi:10.4064/fm-20-1-177-190}.

\bibitem{Lester}
C.~G.~Lester, \emph{Chiral Measurements} (2021),
\href{https://arxiv.org/abs/2111.00623v4}{arXiv:2111.00623v4},
\S3.3, Lemma 3.8.

\bibitem{Fowler}
P.~W.~Fowler, \emph{Quantification of chirality: attempting the impossible},
Symmetry: Culture and Science \textbf{16}(4) (2005), 321--334.
\href{https://journal-scs.symmetry.hu/wp-content/uploads/2014/01/fowler-05-4.pdf}{Full text}.

\bibitem{Wang}
H.~Wang, \emph{Two-component Pseudovectoral Chirality Function for
Tetrahedra} (2021),
\href{https://arxiv.org/abs/2103.12949v1}{arXiv:2103.12949v1}.
\end{thebibliography}
\end{document}
